% Thesis Abstract

\cleardoublepage
% \MARKBOTH{Abstract}{Abstract}

\chapter*{Abstract}

% \begin{abstract}
  % \large \textbf{\ThesisTitleFront}
  % \large \textbf{\ThesisTitlePlaintext}
  \large \textbf{\ThesisTitle}

  \vspace{0.5cm}

  \normalsize
  \noindent Monsoon Depressions (MDs) are important synoptic-scale
  systems for India's seasonal rainfall. While most MDs weaken
  rapidly upon landfall due to the loss of oceanic energy and
  increased friction from land interactions, a subset exhibits
  anomalous inland persistence. This study investigates the
  land-surface feedbacks governing this sustainability, focusing on
  the Brown Ocean Effect (BOE), antecedent soil moisture, and the
  orographic influence of the Eastern Ghats.

  \vspace{0.2cm}

  \noindent ERA5 reanalysis data (0.25\(^\circ\) $\times$
  0.25\(^\circ\), hourly) are used as the primary dataset, and
  detected systems are cross-validated against India Meteorological
  Department (IMD) monsoon seasonal reports. MDs are identified using
  850\,hPa relative vorticity thresholds, and their life cycles are
  characterised across multiple Indian summer monsoon seasons using
  track-following diagnostics. Two complementary analytical
  frameworks constitute the primary methodology. Azimuthal
  Hovm\"{o}ller diagrams, constructed by decomposing the MD domain
  into radial annuli at 25\,km intervals extending to 400\,km from
  the storm centre, reveal the spatio-temporal evolution of heat flux
  (surface LHF and SHF), near-surface soil moisture (ERA5
  \texttt{SWVL1}, 0-7\,cm), total precipitation, and 2-metre
  dew-point temperature relative to the MD centre. Track-following
  time series of these variables, stratified by land and ocean
  phases, are used to quantify thermodynamic transitions across the
  land-sea boundary and along the inland track corridor.

  \vspace{0.2cm}

  \noindent Candidate BOE systems were identified using the Andersen
  and Shepherd (2013) criterion-sustained inland LHF $\geq
  70$\,W\,m$^{-2}$ for $\geq$18 hours and contrasted against
  conventional decaying systems. Antecedent soil moisture was
  assessed using an exponentially weighted index to determine the
  pre-storm ground conditions along future inland tracks.

  \vspace{0.2cm}

  \noindent Results indicate that MDs traversing regions with high
  pre-landfall soil moisture maintain inland Latent Heat Flux (LHF)
  levels comparable to the ocean phase, sustaining a
  precipitation-soil moisture feedback loop consistent with the BOE.
  The Eastern Ghats significantly modulate these systems, reducing
  translational speeds; however, a sufficiently pre-moistened land
  surface is found to facilitate speed recovery on the leeward side.
  These findings improve the mechanistic understanding of MD
  longevity, providing critical insights for seasonal flood
  forecasting over the Indian subcontinent.
% \end{abstract}

\vfill
